\section{Infinite families and a comparison panel}
\label{sec:families}

\subsection{Groups of prime exponent}

A single many-to-one power map and its weighted reverse map already create
a noncommutative factor.

\begin{theorem}
\label{thm:prime-exponent-structural}
Suppose \(\exp(G)=p\) is prime and write \(k=|\operatorname{Conj}(G)|\).
Then
\begin{equation}
 \Pram(G)\cong
 \begin{cases}
 M_2(\Q),&k=2,\\
 M_2(\Q)\oplus\Q,&k>2.
 \end{cases}
 \label{eq:prime-exponent-type-structural}
\end{equation}
In the defining representation, the scalar factor in the second line has
multiplicity \(k-2\).
\end{theorem}

\begin{proof}
Let \(e=e_1\in\Q^k\), let \(u=(1,\ldots,1)^{\mathsf T}\), and let
\(c=(c_1,\ldots,c_k)^{\mathsf T}\) be the class-size vector, so \(c_1=1\) and
\(\sum_ic_i=|G|\).  Since \(x^p=1\) for every \(x\in G\), we have
\(\sigma_p(i)=1\) for all \(i\), so \eqref{eq:power-matrix} gives
\(P_p=ue^{\mathsf T}\); and
\(P_p^*=D_G^{-1}P_p^{\mathsf T}D_G
=D_G^{-1}eu^{\mathsf T}D_G=(|G|/c_1)\,e\,(c^{\mathsf T}/|G|)\), so
\begin{equation}
 P_p=ue^{\mathsf T},
 \qquad
 P_p^*=ec^{\mathsf T}.
 \label{eq:prime-rank-one-matrices}
\end{equation}

\emph{The decomposition.}  Put \(W=\Span_\Q\{e,u\}\) and
\[
 U=\{z\in\Q^k:z_1=0,\ c^{\mathsf T}z=0\}.
\]
Since \(\langle z,e\rangle_G=z_1/|G|\) and
\(\langle z,u\rangle_G=c^{\mathsf T}z/|G|\), the space \(U\) is the
class-orthogonal complement of \(W\).  If \(k\geq2\) then \(u\notin\Q e\), so
\(\dim W=2\); the two functionals \(z\mapsto z_1\) and
\(z\mapsto c^{\mathsf T}z\) are linearly independent because
\(c\neq c_1e=e\) when \(k\geq2\), so \(\dim U=k-2\).  Also
\(W\cap U=0\): if \(\alpha e+\beta u\in U\) then its first coordinate gives
\(\alpha+\beta=0\), while
\(c^{\mathsf T}(\alpha e+\beta u)=\alpha+\beta|G|=0\); subtracting and using
\(|G|>1\) forces \(\alpha=\beta=0\).  Hence
\(\Q^k=W\oplus U\) orthogonally.

Both generators vanish on \(U\): for \(z\in U\),
\(P_pz=u(e^{\mathsf T}z)=z_1u=0\) and
\(P_p^*z=e(c^{\mathsf T}z)=0\).  Both preserve \(W\), since
\(P_pe=u\), \(P_pu=u\), \(P_p^*e=c_1e=e\) and \(P_p^*u=(c^{\mathsf T}u)e=|G|e\).
So \(W\) and \(U\) are reducing for \(P_p\) and \(P_p^*\).

\emph{The action on \(W\).}  In the ordered basis \((e,u)\) the four
operators \(P_p,P_p^*,P_pP_p^*,P_p^*P_p\) restrict to
\begin{equation}
 \begin{aligned}
 P_p|_W&=
 \begin{pmatrix}0&0\\1&1\end{pmatrix},
 & P_p^*|_W&=
 \begin{pmatrix}1&|G|\\0&0\end{pmatrix},\\
 (P_pP_p^*)|_W&=
 \begin{pmatrix}0&0\\1&|G|\end{pmatrix},
 & (P_p^*P_p)|_W&=
 \begin{pmatrix}|G|&|G|\\0&0\end{pmatrix}.
 \end{aligned}
 \label{eq:prime-two-by-two}
\end{equation}
Reading these as vectors in \(\Q^4\) with coordinates
\((a_{11},a_{12},a_{21},a_{22})\) gives the matrix
\[
 \begin{pmatrix}
 0&0&1&1\\
 1&|G|&0&0\\
 0&0&1&|G|\\
 |G|&|G|&0&0
 \end{pmatrix},
\]
which is block anti-diagonal and has determinant
\(\pm|G|\,(|G|-1)^2\neq0\) because \(|G|>1\).  Hence the
four restrictions are linearly independent and span \(M_2(\Q)\).

\emph{Conclusion.}  Every word in \(P_p,P_p^*\) of length at least one
vanishes on \(U\), by the previous paragraph, and the span of such words
restricted to \(W\) is all of \(M_2(\Q)\) by the determinant computation.
Therefore
\[
 \PE(G)
 =\{a\oplus\lambda I_U:a\in\End_\Q(W),\ \lambda\in\Q\},
\]
the summand \(\lambda I_U\) arising only from the adjoined identity.  If
\(U=0\), that is \(k=2\), this is \(M_2(\Q)\).  If \(U\neq0\), the element
\(I_W\oplus0\) lies in the algebra, so the algebra is
\(M_2(\Q)\oplus\Q\), and the scalar factor acts on \(U\), which has
dimension \(k-2\).
\end{proof}

The theorem gives a family independent of the
classification of the underlying groups.  It applies to every elementary
abelian \(p\)-group and, for odd \(p\), to every extraspecial group of
exponent \(p\).  All such groups with more than two classes have abstract
algebra \(M_2(\Q)\oplus\Q\), while the multiplicity \(k(G)-2\) of the scalar
factor retains the class count at represented level.

\subsection{Cyclic groups of prime-power order}

There is a complete formula for cyclic prime powers.

\begin{lemma}
\label{lem:tree-chain-decomposition}
Let \(F\) be the linear map on \(\Q^{C_{p^a}}\) defined on coordinate
vectors by \(F\delta_x=\delta_{px}\).  There is an orthogonal rational
decomposition into reducing subspaces for \(F,F^*\) consisting of one chain
of dimension \(a+1\), together with chains of dimension \(j\) having
multiplicity
\[
 m_a=p-2,\qquad
 m_j=p^{a-j-1}(p-1)^2\quad(1\leq j<a).
 \tag{C}\label{eq:tree-chain-multiplicities}
\]
Terms of multiplicity zero are absent.  On a chain of dimension \(j\), the
algebra generated by \(F,F^*\) is \(M_j(\Q)\), and two chains of the same
length are equivalent as modules over it.
\end{lemma}

\begin{proof}
Write \(C_{p^a}=\Z/p^a\) additively and let
\(\ell(x)=a-v_p(x)\) be the number of steps from \(x\) to \(0\) under
\(x\mapsto px\), so that \(\ell(x)=j\) if and only if \(x\) has order
\(p^j\).  There are \(p^{j-1}(p-1)\) vertices at level \(j\geq1\) and one at
level \(0\).  Every vertex at level \(j\leq a-1\) has \(p\)
predecessors, all at level \(j+1\); the predecessors of \(0\) are \(0\)
itself together with the \(p-1\) elements of order \(p\).  Vertices at
level \(a\) are the leaves, and \(F^*\) annihilates them.

Note first that \(F^*\delta_x=\sum_{py=x}\delta_y\), whence
\begin{equation}
 FF^*\delta_x
 =\sum_{py=x}\delta_{py}
 =p\,\delta_x
 \qquad(\ell(x)\leq a-1),
 \label{eq:ff-star-scalar}
\end{equation}
and \(FF^*\delta_x=0\) on the leaves.  Thus \(FF^*=p\cdot I\) off the
leaves.  This is the source of the uniform chain weights below.

Nonradial chains.  Fix a vertex \(v\) with \(1\leq\ell(v)=j\leq a-1\)
and let \(w_0=\sum_yc_y\delta_y\) run over the \((p-1)\)-dimensional space
of sibling differences on the predecessor set of \(v\), so
\(\sum_yc_y=0\).  Put \(w_k=(F^*)^kw_0\) and set \(i=a-j\).  Then \(Fw_0=(\sum_yc_y)\delta_v=0\),
while \eqref{eq:ff-star-scalar} gives
\begin{equation}
 Fw_k=FF^*w_{k-1}=p\,w_{k-1}
 \qquad(1\leq k\leq i-1),
 \label{eq:chain-weights}
\end{equation}
because \(w_{k-1}\) is supported at level \(j+k\leq a-1\).  The vector
\(w_k\) is supported at level \(j+1+k\), so the chain terminates at
\(w_{i-1}\), and \(F^*w_{i-1}=0\) since that vector is
supported on leaves.  Vectors supported at distinct levels are orthogonal
for the uniform form, so the \(w_k\) are an orthogonal rational basis of a
reducing subspace of dimension \(i\).  Counting: level \(j=a-i\) carries
\(p^{a-i-1}(p-1)\) vertices for \(j\geq1\), each contributing \(p-1\)
independent differences, which gives the multiplicity
\(p^{a-i-1}(p-1)^2\) of \eqref{eq:tree-chain-multiplicities}.

Distinct chains are orthogonal.  This does not follow from level
bookkeeping alone, because two chains born at different levels do have
vectors at a common level, and their supports can even be nested.  We verify
it directly.  Since \((F^*)^k\delta_y=\sum_{z\in F^{-k}(y)}\delta_z\), the
chain vector born at \(v\) is
\begin{equation}
 w_k^{(v)}=\sum_{y\in F^{-1}(v)}c_y\,\mathbf1_{F^{-k}(y)},
 \qquad \sum_{y}c_y=0 .
 \label{eq:chain-explicit}
\end{equation}
Let \(v,v'\) be birth vertices at levels \(j\geq j'\) and let
\(w_k^{(v)},w_{k'}^{(v')}\) both be supported at level \(L\), so
\(k=L-j-1\) and \(k'=L-j'-1\), whence \(k'\geq k\).  If \(j=j'\) and
\(v\neq v'\) the fibres \(F^{-1}(v)\) and \(F^{-1}(v')\) are disjoint, hence
so are the sets \(F^{-k}(y)\) occurring in \eqref{eq:chain-explicit}, and the
inner product vanishes.  If \(j>j'\), then \(k'>k\), and for
\(y\in F^{-1}(v)\), \(y'\in F^{-1}(v')\) the intersection
\(F^{-k}(y)\cap F^{-k'}(y')\) is nonempty only if
\(y'=F^{k'-k}(y)\), in which case \(F^{-k}(y)\subseteq F^{-k'}(y')\) and the
intersection has \(p^{k}\) elements.  But every \(y\in F^{-1}(v)\) satisfies
\(F(y)=v\), so for \(k'-k\geq1\) the value
\(F^{k'-k}(y)=F^{k'-k-1}(v)=:y_0'\) is the same vertex for all such
\(y\).  Therefore
\[
 \langle w_k^{(v)},w_{k'}^{(v')}\rangle
 =p^{k}\,c'_{y_0'}\sum_{y\in F^{-1}(v)}c_y=0 .
\]
The same computation against the radial vectors \(r_L=\mathbf1_{\ell^{-1}(L)}\)
gives \(\langle w_k^{(v)},r_L\rangle=p^{k}\sum_yc_y=0\).  Hence all the chains
listed here and below are mutually orthogonal.

The radial chain.  Put \(r_k=\sum_{\ell(x)=k}\delta_x\) for
\(0\leq k\leq a\).  Then
\begin{equation}
 \begin{aligned}
 Fr_0&=r_0, & Fr_1&=(p-1)r_0, & Fr_k&=p\,r_{k-1}\ (k\geq2),\\
 F^*r_0&=r_0+r_1, & F^*r_k&=r_{k+1}\ (1\leq k\leq a-1), & F^*r_a&=0.
 \end{aligned}
 \label{eq:radial-relations}
\end{equation}
These \(a+1\) orthogonal vectors span a reducing subspace.  The predecessor
set of \(0\) contains \(0\), so \(F\) has the eigenvector \(r_0\) and is
not nilpotent here; this is the one chain on which the endpoint
argument for a nilpotent shift is unavailable.

Chains born at the origin.  The vertex \(0\) is the one birth vertex
whose predecessor set is not contained in a single level: \(F^{-1}(0)\)
consists of \(0\) itself together with the \(p-1\) elements of order \(p\).
Among the \((p-1)\)-dimensional space of sibling differences on \(F^{-1}(0)\),
those not involving \(\delta_0\), that is, supported on the level-one
vertices alone, form a subspace of dimension \((p-1)-1=p-2\).  For such a
\(w_0\) we again have \(Fw_0=(\sum_yc_y)\delta_0=0\), and \(w_0\) is supported
at level one, so the chain \(w_k=(F^*)^kw_0\) terminates at \(w_{a-1}\) and has
dimension \(a\).  This gives the multiplicity \(m_a=p-2\), absent when
\(p=2\).  The one remaining direction on \(F^{-1}(0)\), the one involving
\(\delta_0\), is already accounted for by the radial chain, as the dimension
count in \cref{thm:cyclic-prime-powers} confirms.  Note that \(i=a\) does not
arise from the nonradial analysis above, which required \(1\leq j\leq a-1\), so
the length-\(a\) chains are these \(p-2\) chains.

The algebra restricted to each chain is a full matrix algebra.  Let \(W\) be any of the
above chains, with ordered spanning vectors \(w_0,\ldots,w_{i-1}\) in the
nonradial case or \(r_0,\ldots,r_a\) in the radial case, and suppose
\(S\in\End_\Q(W)\) commutes with \(F|_W\) and \(F^*|_W\).  By
\eqref{eq:chain-weights} and \eqref{eq:radial-relations}, \(F^*|_W\) has
one-dimensional kernel, spanned by the top vector \(w_{i-1}\), respectively
\(r_a\).  Hence \(S\) preserves that line and acts on it by some
\(\lambda\in\Q\).  Applying \(F\) repeatedly and using that every weight in
\eqref{eq:chain-weights} and \eqref{eq:radial-relations} is nonzero, the
weights being \(p\), and \(p-1\geq1\) at the last radial step, gives
\(S=\lambda\) on each successive vector in turn.  Therefore
\((\PE|_W)'=\Q\), and
\cref{thm:rational-bicommutant-structural} gives
\[
 \PE|_W=(\PE|_W)''=\End_\Q(W)\cong M_{\dim W}(\Q).
\]
This argument is uniform in the two cases and does not use nilpotency.

Finally, two chains of equal length are isomorphic as modules.  For the
nonradial chains, including those born at the origin,
\eqref{eq:chain-weights} shows that \(F\) acts as \(p\) times the downward
shift and \(F^*\) as the upward shift, with the same weight \(p\) at every
step, independently of length or birth vertex; so any two such chains of the
same length are carried to one another by the linear map matching their
spanning vectors in order.  The radial chain is not of this form, since its
last step has weight \(p-1\) and \(F\) is not nilpotent on it, but it has
dimension \(a+1\), whereas every nonradial chain has dimension at most \(a\),
so no comparison between the two types is required.
\end{proof}

\begin{theorem}
\label{thm:cyclic-prime-powers}
Let \(a\geq1\).  Then
\begin{align}
 \Pram(C_{p^a})
 &\cong
 M_{a+1}(\Q)\oplus\bigoplus_{j=1}^{a}M_j(\Q),
 &&p\text{ odd},                                      \label{eq:cyclic-odd}\\
 \Pram(C_{2^a})
 &\cong
 M_{a+1}(\Q)\oplus\bigoplus_{j=1}^{a-1}M_j(\Q).
 &&                                                   \label{eq:cyclic-two}
\end{align}
Empty sums are omitted.  The \(M_{a+1}\)-module occurs once.  For odd \(p\),
the multiplicity of the \(M_a\)-module is \(p-2\), and for \(1\leq j<a\)
the multiplicity of the \(M_j\)-module is
\[
 m_j=p^{a-j-1}(p-1)^2.
\]
For \(p=2\), the \(M_a\)-block is absent and the same formula gives
\(m_j=2^{a-j-1}\) for \(j<a\).
\end{theorem}

\begin{proof}
Write \(C_{p^a}\) additively.  Its class measure is uniform, so the sole
canonical generator is pullback \(P\) along \(x\mapsto px\), together with
the ordinary transpose \(P^{\mathsf T}\).  The transpose is the map \(F\)
of \cref{lem:tree-chain-decomposition}.  That lemma gives the displayed
simple factors and their multiplicities.  The dimension identity
\[
 p^a=(a+1)+a(p-2)+
 \sum_{j=1}^{a-1}j\,p^{a-j-1}(p-1)^2
\]
shows that the listed chains exhaust the coordinate space.
\end{proof}

\begin{corollary}
\label{cor:cyclic-prime-power-dimensions}
For odd \(p\),
\[
 \dim_\Q\Pram(C_{p^a})
 =(a+1)^2+\sum_{j=1}^{a}j^2,
\]
whereas for \(p=2\) the final summand \(a^2\) is omitted.  These dimensions
depend only on \(a\) and on whether \(p=2\), while the represented
multiplicities retain \(p\).
\end{corollary}

\subsection{Dihedral and quaternion groups}
\label{sec:family-ladder}

Use the presentations
\[
 D_{2n}=\langle a,b:a^{n}=b^2=1,\ bab^{-1}=a^{-1}\rangle,
 \qquad
 Q_{4m}=\langle a,b:a^{2m}=1,\ b^2=a^m,\ bab^{-1}=a^{-1}\rangle,
\]
with all prime-power maps dividing the exponent.  Modular row reduction
followed by rational reconstruction gives the following table.

\begin{proposition}
\label{prop:family-ladder-table}
For \(3\leq n\leq16\) and \(2\leq m\leq8\), the represented algebras
\(A=\Pram(G)\) have the signatures
\[
\begin{array}{lrr|rrr@{\qquad\qquad}lrr|rrr}
G&|G|&k&\dim A&\dim Z&\dim A'&G&|G|&k&\dim A&\dim Z&\dim A'\\
\cline{1-6}\cline{7-12}
D_6   & 6&3& 9&1&1  & D_{22}&22& 7&13&5&5\\
D_8   & 8&5&17&2&2  & D_{24}&24& 9&51&3&3\\
D_{10}&10&4&10&2&2  & D_{26}&26& 8&14&6&6\\
D_{12}&12&6&26&2&2  & D_{28}&28&10&34&4&4\\
D_{14}&14&5&11&3&3  & D_{30}&30& 9&31&4&4\\
D_{16}&16&7&26&2&5  & D_{32}&32&11&41&3&11\\
D_{18}&18&6&18&3&3  &&&&&&\\
\cline{1-6}\cline{7-12}
Q_8   & 8&5&10&2&5  & Q_{24}&24& 9&51&3&3\\
Q_{12}&12&6&26&2&2  & Q_{28}&28&10&34&4&4\\
Q_{16}&16&7&26&2&5  & Q_{32}&32&11&41&3&11\\
Q_{20}&20&8&30&3&3  &&&&&&\\
\end{array}
\]
Multiplicity freeness holds at every rung except
\(D_{16}\), \(D_{32}\), \(Q_8\), \(Q_{16}\), and \(Q_{32}\), where
\(\dim A'>\dim Z\).
\end{proposition}

\begin{proof}
Rational computation as described above; the multiplicity
statement is \cref{cor:multiplicity-free} applied to the commutant and
center columns.
\end{proof}

\begin{theorem}
\label{thm:dihedral-prime-law}
Let \(p\) be an odd prime, put \(m=(p-1)/2\), and let
\[
 d=\operatorname{ord}_{U_p/\{\pm1\}}(2),\qquad q=m/d.
\]
Then
\[
 \Pram(D_{2p})\otimes_\Q\C
 \cong M_3(\C)\oplus\C^{\,d-1+\epsilon_q},
 \qquad
 \epsilon_q=
 \begin{cases}
 0,&q=1,\\
 1,&q>1.
 \end{cases}
\]
Hence
\[
 \dim_\Q\Pram(D_{2p})=8+d+\epsilon_q.
\]
The class-coordinate representation is multiplicity free if and only if
\(q=1\).
\end{theorem}

\begin{proof}
Let \(R\) be the span of the \(m\) rotation-class coordinates and let
\(\mathbf1_R\) be their sum.  The space
\[
 W=\Span_\C\{\delta_1,\delta_s,\mathbf1_R\}
\]
reduces \(\Psi_2,\Psi_p\) and their class-form adjoints.  Their restrictions
to \(W\), computed with class weights \(1,p,p-1\), generate \(M_3(\C)\).

On \(W^\perp\), the operators \(\Psi_p\) and \(\Psi_p^*\) vanish, while
\(\Psi_2\) is the permutation induced by multiplication by \(2\) on
\((\Z/p\Z)^\times/\{\pm1\}\).  This permutation has \(q\) cycles of length
\(d\).  Its eigenvalues are the \(d\)-th roots of unity.  Removing
\(\mathbf1_R\) removes the eigenvalue \(1\) when \(q=1\); for \(q>1\) its
multiplicity becomes \(q-1\).  Thus the algebra on \(W^\perp\) is
\(\C^{d-1+\epsilon_q}\).  The commutator ideal supplies the \(M_3(\C)\)
summand, so the two restrictions form the displayed direct sum.

For \(q>1\), each nontrivial \(d\)-th root occurs with multiplicity \(q\);
for \(q=1\), all eigenspaces on \(W^\perp\) are one-dimensional.  This gives
the last assertion.
\end{proof}

For \(p=17\), \(d=4\) and \(q=2\), so
\[
 \Pram(D_{34})\otimes_\Q\C\cong M_3(\C)\oplus\C^4,
 \qquad \dim_\Q\Pram(D_{34})=13.
\]

\subsection{Rational classes in \texorpdfstring{\(\mathrm{PSL}(2,q)\)}{PSL(2,q)}}

\begin{theorem}
\label{thm:psl2-rational-classes}
Let \(q=p^f\), let \(G=\mathrm{PSL}(2,q)\), and put
\[
 m_\pm=\frac{q\pm1}{\gcd(2,q-1)} .
\]
\begin{enumerate}
\item Let \(x\in G\) have order \(d\) with \(p\nmid d\).  Then \(d\) divides
\(m_+\) or \(m_-\), and the rational class of \([x]\) has \(\varphi(d)/2\)
elements if \(d>2\) and one element if \(d\leq2\).
\item The rational class of a class of \(p\)-elements has at most two
elements.
\item Let \(r\) be an odd prime.  Then \(G\) has a rational class of size
\(r\) if and only if \(\varphi(d)=2r\) for some \(d>2\) dividing \(m_+\) or
\(m_-\).  Such a \(q\) exists if and only if \(2r+1\) is prime.
\end{enumerate}
\end{theorem}

\begin{proof}
(1)  A nontrivial semisimple element of \(G\) lies in a maximal torus, and the
maximal tori of \(G\) are cyclic of order \(m_+\) or \(m_-\); this gives the
divisibility.  Let \(d>2\), let \(T\) be a maximal torus containing \(x\), and
let \(u\) be a good label, so that \(x^u\in T\) also has order \(d\).  For
\(d>2\) we have \(C_G(x)=C_G(x^u)=T\), so if \(gx^ug^{-1}=x\) then
\(gTg^{-1}=T\), that is, \(g\in N_G(T)\).  Since
\(N_G(T)/T\) has order two and acts on \(T\) by inversion, \(x^u\) is
conjugate to \(x\) if and only if \(u\equiv\pm1\pmod d\).  Thus the stabilizer
of \([x]\) in \(U_N\) is the preimage of \(\{\pm1\}\subseteq(\Z/d)^\times\),
of index \(\varphi(d)/2\).  For \(d\leq2\) every good label acts trivially on
\(\langle x\rangle\).

(2)  If \(q\) is even the \(p\)-elements are the involutions and (1) applies.
If \(q\) is odd, every nontrivial unipotent element of \(\mathrm{SL}(2,q)\)
has order \(p\) and there are two unipotent classes, so \(G\) has two
classes of \(p\)-elements in all.

(3)  By (1) and (2) the sizes of the rational classes are the numbers
\(\varphi(d)/2\) for \(d>2\) dividing \(m_+\) or \(m_-\), together with
values at most two, and \(r\geq3\); this is the stated criterion.  If
\(\varphi(d)=2r\) then \(d\) is \(\ell\) or \(2\ell\) with \(\ell=2r+1\)
prime, except for \(r=3\), where also \(d\in\{9,18\}\); in every case \(2r+1\)
is prime.  Conversely if \(\ell=2r+1\) is prime, choose by Dirichlet a prime
\(q\equiv1\pmod{2\ell}\); then \(\ell\) divides \(m_-=(q-1)/2\) and
\(\varphi(\ell)=2r\).
\end{proof}

\begin{corollary}
\label{cor:psl2-no-seven}
No \(\mathrm{PSL}(2,q)\) has a rational class of size \(7\).
\end{corollary}

\begin{proof}
By \cref{thm:psl2-rational-classes}(3) it suffices that \(14\) is not a value
of \(\varphi\).  If \(\varphi(d)=14\) then \(\varphi(p^a)\mid14\) for every
prime power \(p^a\|d\), so \(\varphi(p^a)\in\{1,2,14\}\).  For odd \(p\),
\(\varphi(p^a)=2\) gives \(p^a=3\), while \(\varphi(p^a)=14\) forces
\(p\mid14\), and \(\varphi(7)=6\), \(\varphi(49)=42\).  Hence every factor
contributes a power of two and \(\varphi(d)\) is a power of two.
\end{proof}

For \(3\leq q\leq32\) the sizes \(\varphi(d)/2\) give odd prime values only
for
\[
 \begin{array}{lll}
 \mathrm{PSL}(2,8):\ d=7,9,\ r=3, &
 \mathrm{PSL}(2,13):\ d=7,\ r=3, &
 \mathrm{PSL}(2,17):\ d=9,\ r=3,\\
 \mathrm{PSL}(2,19):\ d=9,\ r=3, &
 \mathrm{PSL}(2,27):\ d=7,14,\ r=3, &
 \mathrm{PSL}(2,23):\ d=11,\ r=5,\\
 \mathrm{PSL}(2,29):\ d=7,14,\ r=3, &
 \mathrm{PSL}(2,32):\ d=11,\ r=5, &
 \end{array}
\]
and for no other \(q\) in that range; in particular
\(\mathrm{PSL}(2,7),\mathrm{PSL}(2,9),\mathrm{PSL}(2,11),\mathrm{PSL}(2,16),
\mathrm{PSL}(2,25)\) have none.

\subsection{A comparison panel}
\label{sec:census}

\begin{theorem}
\label{thm:comparison-panel}
For each of the following 31 groups the rational algebra
\(A=\Pram(G)\) has the stated dimension, center dimension \(z\), and
represented commutant dimension \(c\); the pairs \((d,m)\) give the complex
blocks \(M_d(\C)\) and their multiplicities in \(V_G\otimes_\Q\C\).

\begingroup
\footnotesize
\normalfont
\setlength{\tabcolsep}{5pt}
\begin{longtable}{@{}lrrrrrrl@{}}
\toprule
$G$ & $|G|$ & $k(G)$ & $\exp G$ & $\dim A$ & $z$ & $c$ & $(d,m)$\\
\midrule
\endfirsthead
\toprule
$G$ & $|G|$ & $k(G)$ & $\exp G$ & $\dim A$ & $z$ & $c$ & $(d,m)$\\
\midrule
\endhead
\multicolumn{8}{@{}l}{\itshape Cyclic}\\
$C_2$ & 2 & 2 & 2 & 4 & 1 & 1 & $(2,1)$\\
$C_3$ & 3 & 3 & 3 & 5 & 2 & 2 & $(2,1),(1,1)$\\
$C_4$ & 4 & 4 & 4 & 10 & 2 & 2 & $(3,1),(1,1)$\\
$C_5$ & 5 & 5 & 5 & 5 & 2 & 10 & $(2,1),(1,3)$\\
$C_6$ & 6 & 6 & 6 & 20 & 2 & 2 & $(4,1),(2,1)$\\
$C_8$ & 8 & 8 & 8 & 21 & 3 & 6 & $(4,1),(2,1),(1,2)$\\
$C_{12}$ & 12 & 12 & 12 & 50 & 4 & 4 & $(6,1),(3,1),(2,1),(1,1)$\\
\addlinespace
\multicolumn{8}{@{}l}{\itshape Elementary abelian}\\
$C_2\times C_2$ & 4 & 4 & 2 & 5 & 2 & 5 & $(2,1),(1,2)$\\
$C_2^3$ & 8 & 8 & 2 & 5 & 2 & 37 & $(2,1),(1,6)$\\
$C_3\times C_3$ & 9 & 9 & 3 & 5 & 2 & 50 & $(2,1),(1,7)$\\
\addlinespace
\multicolumn{8}{@{}l}{\itshape Dihedral}\\
$D_8$ & 8 & 5 & 4 & 17 & 2 & 2 & $(4,1),(1,1)$\\
$D_{10}$ & 10 & 4 & 10 & 10 & 2 & 2 & $(3,1),(1,1)$\\
$D_{12}$ & 12 & 6 & 6 & 26 & 2 & 2 & $(5,1),(1,1)$\\
$D_{16}$ & 16 & 7 & 8 & 26 & 2 & 5 & $(5,1),(1,2)$\\
\addlinespace
\multicolumn{8}{@{}l}{\itshape Quaternion}\\
$Q_8$ & 8 & 5 & 4 & 10 & 2 & 5 & $(3,1),(1,2)$\\
\addlinespace
\multicolumn{8}{@{}l}{\itshape Symmetric}\\
$S_3$ & 6 & 3 & 6 & 9 & 1 & 1 & $(3,1)$\\
$S_4$ & 24 & 5 & 12 & 25 & 1 & 1 & $(5,1)$\\
$S_5$ & 120 & 7 & 60 & 49 & 1 & 1 & $(7,1)$\\
$S_6$ & 720 & 11 & 60 & 59 & 3 & 3 & $(7,1),(3,1),(1,1)$\\
$S_7$ & 5040 & 15 & 420 & 225 & 1 & 1 & $(15,1)$\\
$S_8$ & 40320 & 22 & 840 & 442 & 2 & 2 & $(21,1),(1,1)$\\
\addlinespace
\multicolumn{8}{@{}l}{\itshape Alternating}\\
$A_4$ & 12 & 4 & 6 & 10 & 2 & 2 & $(3,1),(1,1)$\\
$A_5$ & 60 & 5 & 30 & 17 & 2 & 2 & $(4,1),(1,1)$\\
$A_6$ & 360 & 7 & 60 & 27 & 3 & 3 & $(5,1),(1,1),(1,1)$\\
$A_7$ & 2520 & 9 & 420 & 65 & 2 & 2 & $(8,1),(1,1)$\\
$A_8$ & 20160 & 14 & 420 & 146 & 3 & 3 & $(12,1),(1,1),(1,1)$\\
\addlinespace
\multicolumn{8}{@{}l}{\itshape Simple of Lie type}\\
$\mathrm{PSL}(2,7)$ & 168 & 6 & 84 & 26 & 2 & 2 & $(5,1),(1,1)$\\
$\mathrm{PSL}(2,8)$ & 504 & 9 & 126 & 29 & 5 & 5 & $(5,1),(1,1),(1,1),(1,1),(1,1)$\\
$\mathrm{PSL}(2,11)$ & 660 & 8 & 330 & 38 & 3 & 3 & $(6,1),(1,1),(1,1)$\\
\addlinespace
\multicolumn{8}{@{}l}{\itshape Nonabelian of order 27}\\
$H_{27}$ & 27 & 11 & 3 & 5 & 2 & 82 & $(2,1),(1,9)$\\
\addlinespace
\multicolumn{8}{@{}l}{\itshape Frobenius of order 20}\\
$F_{20}$ & 20 & 5 & 20 & 17 & 2 & 2 & $(4,1),(1,1)$\\
\bottomrule
\end{longtable}
\endgroup

\end{theorem}

\begin{proof}
Rational row reduction gives the dimensions of the algebra, center, and
commutant, as in \cref{sec:certification}.  It also enumerates the positive
integral solutions of
\begin{equation}
 \sum_j d_j^2=\dim A,\qquad
 \sum_j m_j^2=\dim A',\qquad
 \sum_jd_jm_j=k(G),\qquad
 \#\{j\}=\dim Z(A),
 \label{eq:dimension-system}
\end{equation}
and for every row the listed solution is unique.  Modular receipts at
\(1{,}000{,}003\) and \(1{,}000{,}033\) reproduce every dimension.
\end{proof}

The groups \(D_8\) and \(Q_8\) both have five classes and identical character
tables, with \(\dim A=17\) and \(\dim A=10\); the invariant is therefore
strictly finer than the character table.  Among the symmetric groups
\(S_3,S_4,S_5\), and \(S_7\) give full matrix algebras, whereas \(S_6\) and
\(S_8\) do not.  The scalar block of \(Q_8\) occurs twice in
class-coordinate space, which is where its commutant differs from that of
\(D_8\).  Both \(H_{27}\) and \(C_3\times C_3\) have \(\dim A=5\), with
\(c=82\) and \(c=50\).
